% Migrated from Appendix B (framework) during the JLEO compression pass.
% Long/multi-step derivations and the full platform-survival detail.
\section{Framework: Full Derivations and Survival Detail}
\label{supp:theory}

This section collects the longer derivations underlying the
evidence-allocation framework of the submitted appendix. The appendix
states each result and gives a short proof; the steps that are routine
or descriptive are reproduced here. Notation follows the appendix:
$T_i$ is firm $i$'s participation count, $W_i$ its number of wins, the
always-loser stratum is $\{i:W_i=0\}$, and the reduced-form types $C$
(cartel-deployed losing role) and $G$ (ordinary losing) are devices for
stating the ranking condition rather than observed classes.

\subsection{Monotone posterior under the Poisson ranking}
\label{supp:theory_monotone}

Conditional on survival and zero wins, take participation counts to
follow the Poisson approximation
$T_i\mid\tau_i,q\sim\text{Poisson}(\lambda_q\tau_i)$ for $q\in\{C,G\}$,
with $\tau_i$ the observed exposure length and $\lambda_C>\lambda_G$ the
surviving-set ordering derived from the exit margin (appendix
Equations for $V_G$, $V_C$, and the selection wedge). By Bayes' rule
the posterior odds of $C$ against $G$ at count $t$ are
\begin{equation}
\label{supp:eq_posterior_odds}
  \frac{\Pr(C\mid T_i=t,W_i=0,\tau_i)}{\Pr(G\mid T_i=t,W_i=0,\tau_i)}
  \;=\;
  \frac{\pi}{1-\pi}\;
  e^{-(\lambda_C-\lambda_G)\tau_i}\,
  \left(\frac{\lambda_C}{\lambda_G}\right)^{t},
\end{equation}
where $\pi$ is the prior on $C$. The first factor and the exponential do
not depend on $t$; the term $(\lambda_C/\lambda_G)^{t}$ is strictly
increasing in $t$ whenever $\lambda_C>\lambda_G$. Hence the posterior
odds, and therefore $\Pr(C\mid T_i=t,W_i=0,\tau_i)$, are non-decreasing
in $t$. This is the monotone-likelihood-ratio property of the Poisson
family \citep{karlin1956theory}: the likelihood ratio of two Poisson
densities is monotone in the count, so any threshold or top-$k$ rule on
$t$ (equivalently on $\log(1+T_i)$, a monotone transform) is admissible
for ranking expected investigative value within the stratum. This
establishes the persistence-ranking proposition of the appendix.

\subsection{Exit margin and the selection wedge}
\label{supp:theory_exit}

The ordering $\lambda_C>\lambda_G$ on the surviving set is obtained
endogenously rather than assumed. An ordinary zero-win firm $G$ pays a
participation cost $c_G>0$ and earns expected gross surplus
$m_G(H_{it})$, weakly declining in its loss history $H_{it}$ under
Bayesian updating; it continues while $V_G(H_{it})\ge 0$ and exits
otherwise. A cartel-deployed losing-role firm $C$ values continuation
through a role value $a_C(H_{it})$ that derives from the role rather
than from its own win prospect, so it does not erode with the loss count
the way $m_G$ does. After sufficiently many zero-win periods
$a_C(H_{it})-c_C>m_G(H_{it})-c_G$, so $C$ remains in the risk set while
$G$ is selectively removed. The surviving conditioning set therefore
over-represents the persistent type, which is exactly the content of
$\lambda_C>\lambda_G$: the rate gap is a \emph{selection} consequence of
differential exit, not an exogenous type effect. The role-value
persistence assumption of the appendix is the minimal condition under
which the wedge holds for histories beyond a loss-count threshold
$\bar L$.

\subsection{Optimal stopping: concavity and comparative statics}
\label{supp:theory_stopping}

The enforcer's evidence-acquisition problem chooses stopping depth $K$
to maximize net recovery
$\Pi(K)=V\!\int_0^{K}\rho(k)\,dk-c\,\Phi(K)$ subject to $\Phi(K)\le B$,
where $\rho(k)$ is the marginal recovery density at rank $k$ (the slope
of the award-recall curve), $\Phi(K)=\int_0^K\phi(k)\,dk$ is cumulative
forensic cost in bid rows, $\phi=\Phi'>0$, $c$ is per-firm forensic
cost, $V$ the recovery payoff per adjudicable case, and $B$ the budget.

The benefit term $V\!\int_0^K\rho$ is concave because $\rho$ is
non-increasing (inherited from the monotone posterior, which makes the
recovery density non-increasing in rank), and the cost term $c\,\Phi$ is
convex because $\phi$ is non-decreasing; their difference $\Pi$ is
concave. The first-order condition $V\rho(K^{*})=c\,\phi(K^{*})$, i.e.
$\rho(K^{*})/\phi(K^{*})=c/V$, is therefore necessary and sufficient for
the interior optimum: the agency descends the ranking until the marginal
case recovered per unit forensic cost falls to the cost--value ratio.
Differentiating the optimality condition in $c/V$ gives
$\mathrm{d}K^{*}/\mathrm{d}(c/V)\le 0$, so cheaper or more valuable
forensics push the agency deeper. Under a binding budget the optimum is
the largest $K$ with $\Phi(K)\le B$, and the multiplier $\mu$ on the
budget plays the role of $c$. Sweeping $c/V$ over $(0,\infty)$ and
reading off $(\Phi(K^{*}),\,\mathrm{recall}(K^{*}))$ traces the
cost--recall frontier reported in the main text.

The proposition is anchored to the \emph{award}-recall curve, where
monotonicity is inherited from the posterior. The two-stage
award-to-bid rerank need not be monotone in the survivor pool: a larger
pool can dilute rerank precision at a fixed depth, and the realized
sequential recall is non-monotone in the data (124 true positives at a
first-stage cut of $1{,}000$ against $116$ at $2{,}000$ for the top
$500$). That is a second-order property of the rerank stage, not of the
award marginal-recovery density, so concavity is claimed for the award
envelope and its comparative statics in $c/V$, not for the height of any
single sequential point.

\subsection{Platform-survival audit: full hazard detail}
\label{supp:theory_survival}

The exit margin predicts who remains on the platform. We examine this
descriptively on a firm-year panel of the $16{,}843$ always-losers
($101{,}366$ firm-years, 2009--2019, right-censored at 2019). Because
the panel uses \emph{full-sample} labels, the exercise is retrospective
and at most mechanism-supporting; ``survival'' means continued BEC
\emph{participation}, not legal-entity survival.

The discrete-time hazard model is a complementary-log-log specification
with duration and calendar-year fixed effects, clustered by firm. Relative
to non-FL always-losers, the exit-hazard ratios are $0.231$ for FL
cobidders and $0.365$ for FL non-cobidders---both far below one, so
exposure to the labels is associated with substantially lower exit
hazard. The cobidder coefficient is sharply estimated
($z=-22.2$, $p<10^{-100}$), and the Kaplan--Meier curves are clearly
separated. The direction is stable across robustness variants:
alternative exit definitions, excluding late cohorts, and excluding the
largest CADE case (ratio $0.281$). A timing-disciplined to-date label,
which restricts the cobidder set to adjudications observable by the
scoring date, attenuates the cobidder ratio to $0.768$ but leaves it
below one. A continuous-time Cox specification yields the same ordering.
Table~\ref{tab:survival_summary_submission} gives the summary moments and
Figure~\ref{fig:survival_curves_submission} the Kaplan--Meier curves. The
audit is descriptive and is not a test of cartel membership.

\begin{table}[htbp]\centering
\begin{adjustbox}{max width=\textwidth,center}
\begin{threeparttable}
\caption{Platform survival of always-losers, by full-sample label.}
\label{tab:survival_summary_submission}\scriptsize
\begin{tabular}{lrrrrrr}
\toprule
Group & $N$ firms & Mean active & Exit share & Median & Right-cens. & Hazard \\
 & & years & $<2019$ & duration & share & ratio \\
\midrule
Non-FL always-loser & $14{,}104$ & $1.61$ & $0.85$ & $1$ & $0.15$ & $1.00$ \\
FL non-cobidder     & $2{,}391$  & $2.87$ & $0.73$ & $2$ & $0.27$ & $0.365$ \\
FL cobidder         & $341$      & $3.51$ & $0.63$ & $3$ & $0.37$ & $0.231$ \\
\bottomrule
\end{tabular}
\begin{tablenotes}\scriptsize
\item Firm-year panel of $16{,}843$ always-losers, $101{,}366$
firm-years, 2009--2019, right-censored at 2019. ``Active years'' counts
distinct years of platform participation; ``duration'' is years between
first and last appearance. Hazard ratios are adjusted discrete-time
(cloglog) estimates relative to the non-FL always-loser baseline. Labels
are full-sample, so the audit is descriptive and mechanism-supporting at
most, not a test of cartel membership.
\end{tablenotes}
\end{threeparttable}
\end{adjustbox}
\end{table}

\begin{figure}[htbp]\centering
\includegraphics[width=0.86\textwidth]{./output/figures/fig_B_survival_curves_frequent_losers}
\caption{Kaplan--Meier platform-survival curves for always-losers, by
full-sample label.}
\label{fig:survival_curves_submission}
\par\medskip
{\footnotesize\textit{Alt text:} Three Kaplan--Meier curves of the share
of firms still participating on the platform against tenure in years. The
FL-cobidder curve lies highest and declines most slowly; the
FL-non-cobidder curve is intermediate; the non-FL always-loser curve
falls steepest and lies lowest, with most firms gone after one year. \par
\textit{Note:} Right-censored at 2019. Labels are full-sample, so the
comparison is descriptive. ``Survival'' is continued BEC platform
participation, not legal-entity survival.}
\end{figure}

\section{Data Construction: Matching, Dedup, and Timing Detail}
\label{supp:data_detail}

This section carries the longer matching and reconciliation detail
behind the validation labels of the submitted Appendix. The Appendix
states the matching rule and the headline counts; the bookkeeping that
makes those counts reproducible is collected here.

\subsection{Defendant-count reconciliation}

The portfolio of $\valDirectCADE$ BEC-active direct defendants and the
$48$ distinct crossmatch CNPJ are not the same object, and the gap is
mechanical. The CADE--BEC crossmatch file holds $48$ distinct CNPJ that
match a BEC participant by exact $14$-digit string equality. Of these,
$\valFunnelDirectActiveFTM$ are active in the firm-tender-map layer used
to build cobidder edges---i.e.\ they appear as a bidder on at least one
tender-item that the participation panel retains---while the remainder
match a registry record but contribute no usable tender-item edge (they
appear only in award rows dropped by the panel's item-level filters, or
only as the unidentified-bidder sentinel). The direct-defendant set of
$\valDirectCADE$ is the subset that anchors at least one retained
tender-item. The three numbers therefore nest as
$\valDirectCADE\to48\to\valFunnelDirectActiveFTM$: the registry match
($48$), the participation-active subset
($\valFunnelDirectActiveFTM$), and the defendants that actually anchor a
loser-side edge ($\valDirectCADE$). No fuzzy or manual step intervenes
at any stage.

\subsection{Dedup and exclusion rules}

Three exclusions fix the cobidder object and are applied in order.
First, the direct defendants are removed from the candidate cobidder
pool, so a firm that is itself an anchor cannot also count as exposure
to an anchor; this keeps the label a measure of the loser-side
\emph{around} the defendants. Second, the unidentified-bidder sentinel
code \texttt{-1}, which BEC writes when a bidder identifier is missing,
is dropped, since it would otherwise collapse many distinct firms into a
single spurious cobidder. Third, a firm that bids in more than one
adjudicated case, or that appears under more than one zero-padded CNPJ
string for the same legal entity, is counted once: cobidder status is a
firm-level indicator, not an edge count, so multiplicity across cases or
CNPJ variants does not inflate the positive count. After these rules the
full $12$-case portfolio yields $\valMainCobidders$ unique always-loser
cobidders.

\subsection{Per-case timing and benchmark membership}

The conservative benchmark and the temporal holdout both turn on
judgment timing, which is observed for $9$ of the $12$ cases; the $3$
undated cases are placed last in any timing ordering and are therefore
excluded from the conservative pre-$\valConservativeCutoff$ window by
construction. Conduct-onset dates, which would be the conceptually
preferred cut, are not in the case record, so judgment dates are the
only timing variable available. Applying the
$\valConservativeCutoff$-12-31 cut to the same cobidder rule leaves
$\valConservativeCases$ cases, $\valConservativeFD$ crossmatch direct
defendants, and $\valConservativeCobidders$ cobidders
($\valConservativeFL$ frequent losers), all strict subsets of the
full-portfolio counts because a narrower case window can only remove
anchors and hence cobidders. The temporal-holdout split scores
participation on $2009$--$2016$ and tests against labels realized in
$2017$--$2019$; firms with no pre-$2017$ history enter as unrankable
entrants ($\valMainEntrantPositives$ of the $\valMainCobidders$
positives), while the $\valStrictRankablePos$ rankable incumbents carry
a pre-window participation score. Because adjudications routinely
post-date conduct spanning several years, both exercises are
retrospective out-of-time discrimination tests rather than real-time
prediction.
